%%=============================================================================
%% CAPITULO 3 -- PROCEDIMENTOS METODOLOGICOS
%%=============================================================================

\section{Procedimentos metodológicos}

A pesquisa adota abordagem quantitativa e caráter aplicado. Modelos de
linguagem classificarão documentos de uma base do setor público que já possui
uma classificação oficial, atribuída por especialistas, que servirá de
referência para medir acertos e erros. Para cada documento serão registrados a
classe proposta pelo modelo e os sinais de incerteza associados a ela.

A avaliação seguirá as medidas próprias da classificação seletiva: para cada
regra de encaminhamento, calcula-se a cobertura, isto é, a fração de casos
aceitos automaticamente, e o risco, isto é, a taxa de erro entre esses casos.
Comparar as curvas risco-cobertura dos diferentes sinais permitirá dizer qual
deles separa melhor os casos que podem dispensar a revisão humana.
